\section{Institutional Setting and Data}
\label{sec:empirical-facts}

\subsection{The fare stack}
A metered yellow-cab fare is a stack, and its components differ in one
respect that matters throughout this paper: who ends up with the money.
Figure~\ref{fig:feestack} sets them side by side.

Two components are \emph{fare revenue the driver keeps}: a rush-hour
surcharge on weekday pickups between 4:00pm and 8:00pm, and an overnight
surcharge from 8:00pm to 6:00am. Both are set by the TLC, and both were
\$1.00 and \$0.50 respectively until December 18, 2022, and \$2.50 and
\$1.00 thereafter \parencite{tlc2022notice}. They are prices, not taxes.

Four components are \emph{collected from the rider and remitted}: a
\$0.50 MTA State Surcharge, a \$1.00 improvement surcharge, a \$2.50 New
York State congestion surcharge on trips south of 96th Street (since
2019), and a \$0.75 MTA congestion toll south of 60th Street (since
January 5, 2025). For these the driver is a collection agent and nets
nothing; on medallion taxis the technology service provider withholds
the congestion surcharge before funds reach the owner.

The distinction organizes the empirical strategy. The within-day clock
designs identify pass-through from the two \emph{fare} surcharges, which
switch at fixed minutes and so deliver sharp, high-powered, repeatedly
replicated variation. The spatial designs identify it from the two
\emph{remitted} congestion fees, which is the fiscal case the paper is
ultimately about but which offers weaker variation. That the two give
the same answer is the point: the payment interface does not distinguish
them, because on the screen they are indistinguishable. Every component
in Figure~\ref{fig:feestack}, whoever receives it, is added to the
metered fare before the payment screen appears, and the percentage
buttons multiply the sum.

Figure~\ref{fig:screens} shows the mechanism as a rider encounters it,
before any econometrics: the same trip four minutes apart, with every
suggested tip higher after 4:00\,PM because a government surcharge has
entered the base.

% Institutional illustration: the fare stack, who receives each component, and
% which components sit inside the tip base. Permitted TikZ use (cf.
% sections/fig_screens.tex); all empirical figures remain StataBE-generated.
% Illustrative trip: weekday 4:30pm pickup, Manhattan south of 60th St, 2025.
% Metered fare is the estimation sample's pre-4pm mean ($17.70).
\begin{figure}[!htbp]
  \centering
  \singlespacing
  \caption{\capttitle{The Fare Stack: Who Receives Each Component, and What
  the Tip Buttons Multiply}}
  \label{fig:feestack}
  \begin{tikzpicture}[
      font=\footnotesize,
      seg/.style={draw=black!55, line width=0.5pt, minimum height=0.85cm},
      drv/.style={seg, fill=statanavy!35},
      gov/.style={seg, fill=stataorange!45},
      lab/.style={font=\scriptsize, align=center},
      lead/.style={draw=black!45, line width=0.35pt}
    ]
    % ---- scale: 0.561 cm per dollar, total $24.95 -> 14.0 cm ----
    \def\y{0}
    % segment: \seg{xstart}{width}{style}
    \node[drv, anchor=south west, minimum width=9.93cm] at (0,\y) {Metered fare \ \$17.70};
    \node[drv, anchor=south west, minimum width=1.40cm] at (9.93,\y) {};
    \node[gov, anchor=south west, minimum width=0.28cm] at (11.33,\y) {};
    \node[gov, anchor=south west, minimum width=0.56cm] at (11.61,\y) {};
    \node[gov, anchor=south west, minimum width=1.40cm] at (12.17,\y) {};
    \node[gov, anchor=south west, minimum width=0.42cm] at (13.57,\y) {};

    % ---- brace ABOVE: the tip base ----
    \draw[decorate, decoration={brace, amplitude=6pt}, stataorange, line width=0.9pt]
      (0,0.95) -- (13.99,0.95);
    \node[font=\small\bfseries, text=stataorange] at (7.0,1.62)
      {The 20/25/30\,\% buttons multiply this entire width: \$24.95};

    % ---- leader lines + labels BELOW for the narrow segments ----
    \draw[lead] (10.63,0) -- (10.63,-0.75) -- (10.20,-0.75);
    \node[lab, anchor=east] at (10.18,-0.75) {Rush-hour surcharge \$2.50};
    \draw[lead] (11.47,0) -- (11.47,-1.30) -- (11.10,-1.30);
    \node[lab, anchor=east] at (11.08,-1.30) {MTA State Surcharge \$0.50};
    \draw[lead] (11.89,0) -- (11.89,-1.85) -- (11.52,-1.85);
    \node[lab, anchor=east] at (11.50,-1.85) {Improvement Surcharge \$1.00};
    \draw[lead] (12.87,0) -- (12.87,-2.40) -- (11.94,-2.40);
    \node[lab, anchor=east] at (11.92,-2.40) {NYS Congestion Surcharge \$2.50};
    \draw[lead] (13.78,0) -- (13.78,-2.95) -- (12.36,-2.95);
    \node[lab, anchor=east] at (12.34,-2.95) {MTA CBD toll \$0.75};

    % ---- braces BELOW: who keeps it ----
    \draw[decorate, decoration={brace, mirror, amplitude=5pt}, statanavy, line width=0.8pt]
      (0,-3.45) -- (11.33,-3.45);
    \node[font=\scriptsize\bfseries, text=statanavy] at (5.66,-3.95)
      {Driver keeps \$20.20};
    \draw[decorate, decoration={brace, mirror, amplitude=5pt}, stataorange, line width=0.8pt]
      (11.35,-3.45) -- (13.99,-3.45);
    \node[font=\scriptsize\bfseries, text=stataorange, align=center] at (12.67,-4.10)
      {Remitted to government\\[-1pt]\$4.75};
  \end{tikzpicture}

  \vspace{0.9em}

  {\footnotesize
  \begin{tabular}{@{}llll@{}}
    \toprule
    Component & Amount & When it applies & Who receives it \\
    \midrule
    Metered fare            & varies & always              & \textbf{Driver} \\
    Rush-hour surcharge     & \$2.50 & weekdays 4--8pm     & \textbf{Driver} \\
    Overnight surcharge     & \$1.00 & 8pm--6am            & \textbf{Driver} \\
    \addlinespace
    MTA State Surcharge     & \$0.50 & always              & MTA \\
    Improvement Surcharge   & \$1.00 & always              & Taxi improvement fund \\
    NYS Congestion Surcharge& \$2.50 & south of 96th St    & New York State \\
    MTA CBD toll            & \$0.75 & south of 60th St    & MTA \\
    Bridge and tunnel tolls & varies & by route            & Bridge authority \\
    \bottomrule
  \end{tabular}}

  \floatnotes{Illustrative weekday 4:30pm pickup south of 60th Street in 2025;
  the metered fare is the estimation sample's pre-4pm mean. Segment widths are
  proportional to dollars. Navy components are fare revenue the driver keeps;
  orange components are collected from the rider and remitted---the driver is
  a collection agent for them and nets nothing. The rush-hour and overnight
  surcharges are set by the TLC but are \emph{not} government revenue. Every
  component, whoever receives it, sits inside the base the percentage buttons
  multiply, which is the fact this paper measures: the \$4.75 of remitted fees
  on this trip raises the suggested tip by \$0.95 at the 20 percent button, and
  raises the tip actually paid by about \$0.52 at our estimated pass-through---
  money that reaches the driver rather than the agency levying the fee.}
\end{figure}


% Illustrative diagram: how a surcharge enters the tip base.
% Institutional illustration only (permitted TikZ use); all empirical
% figures remain StataBE-generated. Amounts use the estimation sample's
% pre-4pm mean fare ($17.70) and the statutory fare stack.
\begin{figure}[!htbp]
  \centering
  \singlespacing
  \caption{\capttitle{The Mechanism on the Screen: One Trip, Four Minutes Apart}}
  \label{fig:screens}
  \begin{tikzpicture}[
      font=\footnotesize,
      screen/.style={rounded corners=8pt, draw=black!55, line width=0.7pt,
        fill=black!2, minimum width=5.35cm, minimum height=6.15cm},
      hdr/.style={font=\footnotesize\bfseries},
      rline/.style={font=\scriptsize},
      btn/.style={rounded corners=3pt, draw=statanavy, fill=statanavy!12,
        minimum width=1.45cm, minimum height=0.86cm, align=center,
        font=\scriptsize},
      btnhot/.style={btn, draw=stataorange, fill=stataorange!18},
      note/.style={font=\scriptsize\itshape, align=center}
    ]
    % ---------------- left screen: 3:58 PM ----------------
    \node[screen] (L) at (0,0) {};
    \node[hdr] at ($(L.north)+(0,-0.42)$) {Pickup 3:58\,PM};
    \node[rline, anchor=west] at ($(L.north west)+(0.38,-1.05)$) {Metered fare};
    \node[rline, anchor=east] at ($(L.north east)+(-0.38,-1.05)$) {\$17.70};
    \node[rline, anchor=west] at ($(L.north west)+(0.38,-1.50)$) {MTA + improvement};
    \node[rline, anchor=east] at ($(L.north east)+(-0.38,-1.50)$) {\$1.50};
    \node[rline, anchor=west] at ($(L.north west)+(0.38,-1.95)$) {Congestion surcharge};
    \node[rline, anchor=east] at ($(L.north east)+(-0.38,-1.95)$) {\$2.50};
    \node[rline, anchor=west, text=black!45] at ($(L.north west)+(0.38,-2.40)$) {Rush-hour surcharge};
    \node[rline, anchor=east, text=black!45] at ($(L.north east)+(-0.38,-2.40)$) {---};
    \draw[black!45] ($(L.north west)+(0.38,-2.72)$) -- ($(L.north east)+(-0.38,-2.72)$);
    \node[rline, anchor=west, font=\scriptsize\bfseries] at ($(L.north west)+(0.38,-3.05)$) {Total};
    \node[rline, anchor=east, font=\scriptsize\bfseries] at ($(L.north east)+(-0.38,-3.05)$) {\$21.70};
    \node[note, text=black!55] at ($(L.north)+(0,-3.72)$) {Add a tip?};
    \node[btn] (L20) at ($(L.north)+(-1.72,-4.45)$) {20\%\\[-1pt]\textbf{\$4.34}};
    \node[btn] (L25) at ($(L.north)+(0,-4.45)$) {25\%\\[-1pt]\textbf{\$5.43}};
    \node[btn] (L30) at ($(L.north)+(1.72,-4.45)$) {30\%\\[-1pt]\textbf{\$6.51}};
    \node[note, text=black!55] at ($(L.south)+(0,0.5)$) {No tip \quad Custom};
    % ---------------- right screen: 4:02 PM ----------------
    \node[screen] (R) at (7.6,0) {};
    \node[hdr] at ($(R.north)+(0,-0.42)$) {Pickup 4:02\,PM};
    \node[rline, anchor=west] at ($(R.north west)+(0.38,-1.05)$) {Metered fare};
    \node[rline, anchor=east] at ($(R.north east)+(-0.38,-1.05)$) {\$17.70};
    \node[rline, anchor=west] at ($(R.north west)+(0.38,-1.50)$) {MTA + improvement};
    \node[rline, anchor=east] at ($(R.north east)+(-0.38,-1.50)$) {\$1.50};
    \node[rline, anchor=west] at ($(R.north west)+(0.38,-1.95)$) {Congestion surcharge};
    \node[rline, anchor=east] at ($(R.north east)+(-0.38,-1.95)$) {\$2.50};
    \node[rline, anchor=west, text=stataorange, font=\scriptsize\bfseries] at ($(R.north west)+(0.38,-2.40)$) {Rush-hour surcharge};
    \node[rline, anchor=east, text=stataorange, font=\scriptsize\bfseries] at ($(R.north east)+(-0.38,-2.40)$) {\$2.50};
    \draw[black!45] ($(R.north west)+(0.38,-2.72)$) -- ($(R.north east)+(-0.38,-2.72)$);
    \node[rline, anchor=west, font=\scriptsize\bfseries] at ($(R.north west)+(0.38,-3.05)$) {Total};
    \node[rline, anchor=east, font=\scriptsize\bfseries] at ($(R.north east)+(-0.38,-3.05)$) {\$24.20};
    \node[note, text=black!55] at ($(R.north)+(0,-3.72)$) {Add a tip?};
    \node[btnhot] (R20) at ($(R.north)+(-1.72,-4.45)$) {20\%\\[-1pt]\textbf{\$4.84}};
    \node[btnhot] (R25) at ($(R.north)+(0,-4.45)$) {25\%\\[-1pt]\textbf{\$6.05}};
    \node[btnhot] (R30) at ($(R.north)+(1.72,-4.45)$) {30\%\\[-1pt]\textbf{\$7.26}};
    \node[note, text=black!55] at ($(R.south)+(0,0.5)$) {No tip \quad Custom};
    % ---------------- annotations ----------------
    \draw[-{Stealth[length=2.6mm]}, line width=0.9pt, stataorange]
      ($(L.east)+(0.12,0.35)$) -- ($(R.west)+(-0.12,0.35)$)
      node[midway, above, note, text=stataorange, text width=1.9cm]
      {4:00\,PM passes};
    \draw[decorate, decoration={brace, mirror, amplitude=5pt}, stataorange, line width=0.8pt]
      ($(R20.south west)+(-0.06,-1.75)$) -- ($(R30.south east)+(0.06,-1.75)$);
    \node[note, text=stataorange, text width=4.6cm] at ($(R.south)+(0,-0.5)$)
      {each button rises by its percentage of the \$2.50 surcharge};
  \end{tikzpicture}
  \floatnotes{Illustration of the identifying variation, using the
  estimation sample's pre-4pm mean metered fare. The two screens show the
  same trip picked up four minutes apart on a weekday. At 4:00\,PM the
  \$2.50 rush-hour surcharge (in orange) enters the total, and because
  the suggested tips are percentages of that surcharge-inclusive total,
  every button rises with it---the 20 percent button by \$0.50. Nothing
  about the ride, the driver, or the rider's intended generosity has
  changed. The paper measures how much of the surcharge reaches the
  driver through this channel: about fifteen cents per surcharge dollar,
  averaged over riders who press a button, tip a custom amount, or skip
  the tip.}
\end{figure}


\subsection{What base do the buttons use?}
Figure~\ref{fig:spikes} shows the answer before any matching exercise.
Expressed as a percentage of the all-in pre-tip total, 57.3 percent of
positive credit-card tips fall in the half-point bin at \emph{exactly}
20 percent, with further spikes at exactly 25 and 30---the three
buttons. Expressed as a percentage of the metered fare alone, the same
tips have no spikes anywhere: no bin holds more than 3.6 percent. The
smaller spikes at 10 and 15 percent of the all-in total are custom
round percentages computed on the same inclusive base---a first sign,
confirmed in Section~\ref{sec:results-vendor}, that the base governs
even riders who never press a button.

\begin{figure}[!htbp]
  \centering
  \singlespacing
  \caption{\capttitle{The Base, Drawn: The Same Tips Against Two
  Candidate Bases}}
  \label{fig:spikes}
  \includegraphics[width=0.88\textwidth]{fig-spikes.pdf}
  \floatnotes{June 2025, standard-rate credit-card trips with positive
  tips and no tolls (2.18 million trips), in bins of half a percentage
  point over 5--40 percent. Panel A expresses each tip as a percentage
  of the all-in pre-tip total---the base the buttons multiply;
  annotations mark the bins at exactly 20, 25, and 30 percent. Panel B
  expresses the same tips as a percentage of the metered fare alone, on
  identical axes. Only the buttons' base produces spikes.}
\end{figure}

Table~\ref{tab:vendor} sharpens the fact by vendor. For each
payment vendor we compute the share of credit-card tips that equal 20,
25, or 30 percent of a candidate base to within half a cent. In July
2024, before the CBD toll existed, the fee-inclusive and fee-exclusive
bases are numerically identical and the match rates coincide exactly, as
they must. In June 2025 they diverge sharply: the fee-inclusive total
matches 17 to 74 percent of tips depending on vendor, the fee-exclusive
total only 4 to 17 percent. All three vendors compute suggested tips on
the surcharge-inclusive total.

The table also documents striking vendor heterogeneity in adherence: at
one vendor 17 percent of tips are exact button presses, at the others 66
and 74 percent. Our estimates aggregate over this heterogeneity, so
$\rho$ should be read as a market-wide average rather than a
vendor-specific behavioral parameter.

\begin{table}[!htbp]
  \centering
  \singlespacing
  \caption{\capttitle{Which Base Do the Percentage Buttons Use?}}
  \label{tab:vendor}
  \small
  \begin{tabular}{@{}lcccc@{}}
    \toprule
    & & \multicolumn{3}{c}{Share of tips matching 20/25/30\% of:} \\
    \cmidrule(lr){3-5}
    Vendor, month & Card tips & All-in total & Excl.\ CBD fee & Fare only \\
    \midrule
    \input output/tables/tab-vendor.tex
    \bottomrule
  \end{tabular}
  \floatnotes{Card trips with positive tips and no tolls. July 2024
  precedes the CBD congestion fee, so the first two bases are identical
  by construction; June 2025 follows it. A match requires equality to
  within half a cent.}
\end{table}

\subsection{Data and estimation cells}
Trip records for every yellow-cab trip---pickup timestamp to the second,
fare components, tip (recorded for card payments only), payment type,
vendor, and pickup zone---are published monthly by the TLC
\parencite{tlc2026trip}. We aggregate to date-by-minute cells in windows
around the three clock cutoffs, recording trip counts, mean metered
surcharge, mean pre-tip total, mean card tip, the share of tips exactly
at a default button, the zero-tip share, and the card share. Separate
zone-by-fare-bin cells support the congestion-fee designs. Cells, not
raw trips, are the estimation input; the build is deterministic and
versioned with the code.

Table~\ref{tab:descriptive} describes the sample. The two surcharge
regimes are similar in composition---card share near 77 percent, mean
tips of \$3.44 and \$4.49, mean pre-tip totals of \$19 and \$26---with
the difference in the mean metered surcharge reflecting the December
2022 increase. Roughly half of all credit-card tips are exactly equal to
a default button, which is the behavioral fact the rest of the paper
builds on.

\begin{table}[!htbp]
  \centering
  \singlespacing
  \caption{\capttitle{Descriptive Statistics by Surcharge Regime}}
  \label{tab:descriptive}
  \small
  \begin{tabular}{@{}p{0.62\textwidth}r@{}}
    \toprule
    & Value \\
    \midrule
    \input output/tables/tab-descriptive.tex
    \bottomrule
  \end{tabular}
  \floatnotes{Date-by-minute cells in the estimation windows around the
  4:00pm, 8:00pm, and 6:00am cutoffs, September 2021--September 2024.
  Trip-weighted means. Tips are observed for credit-card payments only.}
\end{table}

\subsection{The motivating pattern}
Figure~\ref{fig:firststage} shows the treatment. Mean metered surcharge
by minute of day rises in a single step at 4:00pm---by \$0.72 in the
\$1.00 regime and \$1.95 in the \$2.50 regime. The step is sharp because
the surcharge is assessed on the pickup timestamp; were it assessed at
dropoff, the trip-duration distribution would smear it across half an
hour.

The steps fall short of the statutory \$1.00 and \$2.50, and the reason
is a property of the data rather than of the fare rules. The TLC's
\texttt{extra} field records miscellaneous extras, not the rush
surcharge alone, and it is contaminated on both sides of the cutoff. In
the \$2.50 regime, 23.0 percent of standard-rate card trips in the hour
\emph{before} 4:00pm already carry exactly \$2.50 in that field, while
4.5 percent of trips in the hour after carry nothing; the balance sits
at values such as \$5.00 and \$7.50 that no single surcharge explains
(Table~\ref{tab:composition}). Because the measured step is therefore
not a clean statutory dose, we report two scalings throughout: an
intention-to-treat effect per \emph{statutory} dollar, which divides by
a known constant and requires no assumption about the field, and an
instrumental-variables $\rho$ that scales by the measured step. The
first is the conservative number and the second requires the
discontinuity in \texttt{extra} to be entirely rush surcharge---an
assumption the composition table does not establish. We report both and
flag which is which at every point.

\begin{figure}[!htbp]
  \centering
  \singlespacing
  \caption{\capttitle{The First Stage: Metered Surcharge by Minute of Day}}
  \label{fig:firststage}
  \includegraphics[width=0.82\textwidth]{fig-firststage.pdf}
  \floatnotes{Trip-weighted mean of the metered surcharge field by
  minute of pickup, weekdays excluding holidays. Navy: \$1.00 regime
  (September 2021--December 2022). Orange: \$2.50 regime (December
  2022--September 2024). The dashed line marks 4:00pm.}
\end{figure}

\begin{table}[!htbp]
  \centering
  \singlespacing
  \caption{\capttitle{What the Surcharge Field Is Made Of}}
  \label{tab:composition}
  \small
  \begin{tabular}{@{}lccccc@{}}
    \toprule
    & \multicolumn{4}{c}{Share of card trips with \texttt{extra} equal to:}
      & Mean \\
    \cmidrule(lr){2-5}
    Cutoff, side & \$0 & \$1.00 & \$2.50 & other & \texttt{extra} \\
    \midrule
    \input output/tables/tab-composition.tex
    \bottomrule
  \end{tabular}
  \floatnotes{Standard-rate card trips, weekdays excluding holidays,
  \$2.50 regime, within 60 minutes of each cutoff. The field is
  contaminated on both sides of every cutoff: a substantial share of
  trips carries the higher value before the surcharge switches on, and
  a nonzero share carries none after. This is why the measured step
  falls short of the statutory dose, and why estimates scaled by the
  measured step ($\rho^{IV}$) are reported alongside estimates scaled by
  the statutory dose (ITT).}
\end{table}
